% \hyphenchar and \skewchar: per-font integers, readable in the gullet.
%
% The stomach assigned them and threw the value away (`let _c = ...`), and the
% gullet could not read them at all -- `\the\hyphenchar\tenrm` handed back its
% own text. Both halves were wrong in different ways, which is the pattern
% behind most of the faults in this engine.
%
% It is not a niche primitive. expl3 uses \hyphenchar as general-purpose
% per-font storage: l3intarray keeps an integer array's *length* there, and
% reads it back as
%
%   \intarray_count:N #1 -> \int_value:w \__intarray_count:w #1
%
% where \__intarray_count:w is \hyphenchar. Discarding the value lost every
% array's size.
%
% A newly loaded font takes both from \defaulthyphenchar and \defaultskewchar
% as they stand at load time (verified against INITEX: 0, 0, then 88 after
% setting it). plain.tex sets the former to 45, so its fonts hyphenate on `-`.
\input plain
\def\M#1{\message{<#1>}}
\M{a=\the\hyphenchar\tenrm}
\M{b=\the\skewchar\tenrm}
\hyphenchar\tenrm=45
\M{c=\the\hyphenchar\tenrm}
\M{d=\number\hyphenchar\tenrm}
\M{e=\ifnum\hyphenchar\tenrm=45 T\else F\fi}
\font\seven=cmr7 \M{f=\the\hyphenchar\seven}
\hyphenchar\seven=99 \M{g=\the\hyphenchar\seven}
\M{h=\the\hyphenchar\tenrm}
\skewchar\seven=7 \M{i=\the\skewchar\seven}
\M{k=\the\hyphenchar\font}
\M{l=\ifnum\hyphenchar\seven>0 pos\else nonpos\fi}
\edef\x{\the\hyphenchar\seven}\M{m=\x}
\bye
